\documentclass[10pt,a4paper]{amsart}
\usepackage[margin=19mm]{geometry}
\usepackage{amsmath,amssymb,mathtools}
\usepackage[scr=rsfs,cal=euler]{mathalfa}
\usepackage{xfrac}
\usepackage[unicode,hidelinks,bookmarks=false]{hyperref}
\newcommand{\ie}{i.e.\ }
\newcommand{\cf}{cf.\ }
\newcommand{\resp}{respectively\ }
\newcommand{\Subsection}{\subsection}
\providecommand{\nobreakdash}{\nobreak\hbox{-}}
\setlength{\emergencystretch}{2em}
\multlinegap0pt
\usepackage{bm}
\usepackage{bbm}
\usepackage{dsfont}
\usepackage{xspace}
\usepackage{cancel}
\usepackage{mathtools}
\usepackage{amsthm}
\makeatletter
\@ifundefined{theorem}{\newtheorem{theorem}{Theorem}}{}
\@ifundefined{lemma}{\newtheorem{lemma}[theorem]{Lemma}}{}
\makeatother
\title[A logarithmic approximation of linearly ordered colourings]{A logarithmic approximation of linearly ordered colourings}
\author{Johan H{\aa}stad, Bj\"orn Martinsson, Tamio-Vesa Nakajima, Stanislav \v{Z}ivn\'y}
\date{Source: arXiv:2404.19556v6 (CC BY 4.0)}
\begin{document}
\maketitle

\noindent\emph{Source and licence.} A logarithmic approximation of linearly ordered colourings. Version arXiv:2404.19556v6 (CC BY 4.0), distributed under the Creative Commons Attribution 4.0 International licence, as recorded in the arXiv licence metadata for this exact version and shown on the abstract page https://arxiv.org/abs/2404.19556v6. Full rights notice: licenses/arxiv-2404.19556-CC-BY-4.0.txt.\par\smallskip

\noindent\textit{Editorial scope.} Counted unit: the Lemma whose source label is \texttt{lem:innerstep} in the article (source0.tex lines 151--156, source label \texttt{lem:innerstep}), delivered together with the article's own complete original proof of it (source0.tex lines 158--182, one \texttt{proof} environment of 25 lines). No mathematics is written, repaired or re-derived: statement and proof are the article's own text line for line. Required context retained uncounted: none. Every definition, display and label of those lines is retained unchanged. Omitted from this package and not redistributed: 1--150, 157--157, 183--361. Nothing inside a retained excerpt is omitted or altered: every retained body is delivered exactly as the article's lines are written, so no omission receipt of any kind accompanies this package. Citations: the retained excerpts carry no citation command at all, so no bibliography entry of the article is transcribed and no reference is redistributed. Reference adaptation: none was needed and none was made; every reference inside the delivered unit resolves inside it, so no unresolved reference or citation is produced. Printed slips kept literal: no printed word, symbol, hypothesis or number of the counted statement or of its proof was altered. Layout and class: the article's own class is \texttt{article} with packages including amssymb, amsmath, amsthm, babel, microtype, newtx, a4wide, anyfontsize, complexity, xcolor, hyperref, cleveref; the wrapper is the lane's pinned \texttt{amsart} editorial wrapper at 10pt on A4, which restates the theorem-like environments the retained text needs and adds the article's own macro definitions that the retained text needs (none), with extra wrapper packages (bm, bbm, dsfont, xspace, cancel, mathtools). The delivered numbering is the unit's own. Assets: no retained span contains a figure environment, a graphics-inclusion command, a PDF-page inclusion, a table or a TikZ picture; no image file is copied into this package, the package declares no asset dependency and no image file is redistributed with it. Licence: version arXiv:2404.19556v6 of this article is distributed under the Creative Commons Attribution 4.0 International licence, as recorded in the arXiv licence metadata for this exact version and shown on the abstract page https://arxiv.org/abs/2404.19556v6. A full rights notice is delivered at licenses/arxiv-2404.19556-CC-BY-4.0.txt. Characters: every retained excerpt is pure ASCII, so the delivered engine, XeLaTeX with its default Latin Modern text font, reports no missing character.
\begin{lemma}\label{lem:innerstep}
    There is an algorithm which, if given a $\{2, 3\}$-uniform hypergraph $H$ with $n$ vertices and
    $m$ edges that admits an LO 2-colouring and such that the implied linear
    system of equations $Av=1^m$ does not fix the value of any variable, outputs a subset $T$ of vertices that intersects edges of size three in zero or two vertices and edges of size two in exactly one vertex.
    Moreover, we have  $\vert T\vert \geq n / 2$. The algorithm runs in $O(n^3 + nm)$ time.
\end{lemma}

\begin{proof}

    We first describe a randomised version of our algorithm, and then derandomise
    it. The set of solutions to $Av=1^m$ is an affine space and hence a generic solution can be written as $v=v^0+\sum_{i=1}^r a_i v^i$ for a basic solution $v^0$, linearly independent solutions to the homogeneous system $v^i$, and field elements (in this case bits) $a_i$.
    The fact that no variable is fixed implies that for each vertex $x$ there is 
    some positive $i$ such that $v^i_x=1$.
    
     For the randomised algorithm choose $a_1, \ldots, a_r$ to be independent identically distributed uniformly random bits, and set $T$ to be 
     the set of variables $x$, such that $v_x=0$.  Clearly $T$ satisfies
     the conclusion of the lemma as in each edge we have an odd number of ones. 
     For every vertex $x$, there exists positive $i$ such that $v_x^i = 1$ --- hence, due to the influence of $a_i v_x^i$, we see that $x$ is included in $T$ with probability $\frac{1}{2}$.
     Thus, on average $T$ contains half the vertices.

    Now, we derandomise this algorithm using the method of conditional
    expectations. Go through the variables $a_i$ in increasing order and 
    fix its value once and for all.  Fixing the value of $a_i$ determines
    the value(s) of some $v_x$ while other values remain undetermined.  For each
    value being determined $v_x^i=1$ and hence one value of $a_i$ gives
    the final value 0 and the other gives final value 1.  Set $a_i$ such
    that at least half the determined values are 0.  After we have 
    fixed all $a_i$ this way, we have a final solution with at least
    $n/2$ zeroes.
    
    The bottleneck of the running time of this algorithm is solving the linear system of equations. This can be done in the advertised running time since every equation has $O(1)$ entries. 
\end{proof}

\end{document}
